Proiectul a urmărit construirea unui sistem de cache in-memory distribuit, de la structura de date până la
infrastructura de producție, și evaluarea lui riguroasă. Principalele concluzii:

\begin{enumerate}
  \item \textbf{Cache-ul accelerează considerabil, dar numai după ce baza de date este interogată corect.}
        Profilarea a arătat că gâtul de sticlă inițial era PostgreSQL (410\% CPU). Tuning-ul, indecșii
        acoperitori, eliminarea serializării prin \texttt{NOTIFY} și agregatele materializate au crescut
        debitul cu cache de la \R{opt/baseline/rps} la \R{opt/mat-local/rps} cereri/s și au redus p99 de la
        \R{opt/baseline/p99} la \R{opt/mat-local/p99}~ms. Pe această bază optimizată, cache-ul local mai aduce
        \R{e1/local/speedup}× debit și reduce cu 27\% CPU-ul consumat per cerere.
  \item \textbf{Proximitatea contează.} Cache-ul în proces are cea mai mică latență. Cache-urile partajate
        (clusterul propriu, Redis) au rată de hit mai mare și coerență mai ieftină. Un cache la distanță este
        rentabil doar peste un prag al ratei de hit (E2).
  \item \textbf{O tabelă in-memory într-un SGBD relațional nu este un cache eficient} pentru date pe care SGBD-ul
        le servește deja din memorie: costul unui hit este cel al unei interogări SQL (E1, E8).
  \item \textbf{Strategia de scriere este un compromis consistență–performanță.} Write-behind cu benzi
        partiționate oferă cel mai mare debit și reduce tranzacțiile de câteva ori, cu consistență eventuală.
        Write-through păstrează rata de hit cu consistență imediată (E3).
  \item \textbf{Paralelismul are limite concrete.} Sharding-ul transformă un motor care se blochează (1 shard)
        într-unul care scalează, dar cheile fierbinți limitează accelerarea la acces Zipf. Scalarea orizontală
        aduce beneficii doar cu hardware suplimentar (E5, ENG).
  \item \textbf{Mecanismele de robustețe merită costul}: single-flight reduce cu 29\% încărcarea bazei la
        pornirea la rece, invalidarea prin \texttt{LISTEN/NOTIFY} menține coerența în toate cele 18
        configurații, iar write-behind nu pierde date la oprire.
\end{enumerate}

Toate obiectivele din \cref{ch:intro} au fost atinse: motorul de tabele in-memory, cache-ul distribuit,
aplicația web cu backend-uri și strategii comutabile, comparația cu Redis și cu tabela in-memory relațională
și evaluarea experimentală reproductibilă.
